\section{Scheduling}

\begin{figure}[t!]
  \centering
  \includegraphics[width=0.65\linewidth]{figures/sample_pallet.pdf}
  \caption{Pallet layout produced by the grouping stage: 51 groups
  tiling the $11 \times 17$ grid, each colored
  distinctly. Groups alternate $3$--$4$--$4$ across each row.}
  \label{fig:sample-pallet}
\end{figure}

The second part of the model determines the release order of the groups and
selects which gripper plate opens for each release. In this case, the most
direct representation consists in assigning to each group a single release
position. The array \texttt{drop\_order} stores the time step at which each
group is released, while \texttt{plate} stores the plate opened for that release.
The value \(0\) denotes the left plate, and the value \(1\) denotes the right
plate.

\vspace{0.5cm}
\begin{lstlisting}
array[1..n_groups] of var 1..n_groups: drop_order;
array[1..n_groups] of var 0..1: plate;
\end{lstlisting}
\vspace{0.5cm}

The only global constraint required to make \texttt{drop\_order} a valid
schedule is \texttt{alldifferent}. Since every value in \texttt{drop\_order}
belongs to the domain \(1,\dots,\texttt{n\_groups}\), forcing all values to be
different ensures that each release step is assigned to exactly one group.

\vspace{0.5cm}
\begin{lstlisting}
constraint alldifferent(drop_order);
\end{lstlisting}
\vspace{0.5cm}

The feasibility of a release depends on the position of the adjacent groups in the same
vertical column. As shown in Figure~\ref{fig:sample-pallet}, groups are
numbered from left to right within each row, and then from top to
bottom across rows.
Therefore, since each row contains \texttt{groups\_per\_row} groups,
two groups that are aligned in consecutive rows have indices
that differ by \texttt{groups\_per\_row}. For a generic group \(g\),
the vertical neighborhood is therefore given by 
$g \pm \texttt{groups\_per\_row},$ whenever these indices are within the valid range.
Groups in the first pallet row have no upper neighbor,
while groups in the last pallet row have no lower neighbor.

The final constraint encodes the physical restriction on the gripper plates:
when a group is released, the selected plate must open toward a free direction.
A direction is free if the group lies on the corresponding pallet boundary, or
if the group in that direction is scheduled to be released later and is
therefore still absent from the pallet at the current release step.

Since the relevant neighbors are above and below the current group, the
constraint is written by separating the first row, the internal rows, and the
last row. This avoids invalid array accesses at the pallet boundaries.

\vspace{0.5cm}
\begin{lstlisting}
constraint forall(g in 1..n_groups where g <= groups_per_row)(
  plate[g] = 0 \/
  (
    plate[g] = 1 /\
    drop_order[g + groups_per_row] > drop_order[g]
  )
);

constraint forall(g in 1..n_groups where
  g > groups_per_row /\ g <= n_groups - groups_per_row
)(
  (
    plate[g] = 0 /\
    drop_order[g - groups_per_row] > drop_order[g]
  )
  \/
  (
    plate[g] = 1 /\
    drop_order[g + groups_per_row] > drop_order[g]
  )
);

constraint forall(g in 1..n_groups where g > n_groups - groups_per_row)(
  (
    plate[g] = 0 /\
    drop_order[g - groups_per_row] > drop_order[g]
  )
  \/
  plate[g] = 1
);
\end{lstlisting}
\vspace{0.5cm}

The first constraint handles groups in the first pallet row. These groups have
no upper neighbor, so opening plate \(0\) is always feasible. Opening plate
\(1\), instead, is feasible only if the lower neighbor is released later.

The second constraint handles the internal rows. In this case, both vertical
neighbors exist. Opening plate \(0\) is feasible only if the upper neighbor
\(g-\texttt{groups\_per\_row}\) is released later. Similarly, opening plate
\(1\) is feasible only if the lower neighbor
\(g+\texttt{groups\_per\_row}\) is released later.

The third constraint handles groups in the last pallet row. These groups have
no lower neighbor, so opening plate \(1\) is always feasible. Opening plate
\(0\) is feasible only if the upper neighbor is released later.

The model is formulated as a feasibility problem:

\vspace{0.5cm}
\begin{lstlisting}
solve satisfy;
\end{lstlisting}
\vspace{0.5cm}

No objective function is required. Any assignment of \texttt{drop\_order} and
\texttt{plate} that satisfies the vertical-neighbor constraint defines a valid
pick-and-place plan. Figure~\ref{fig:sample-schedule} shows an example of such a
schedule at different placement steps.

\begin{figure}[h]
  \centering
  \begin{subfigure}[t]{0.24\linewidth}
    \centering
    \includegraphics[width=\linewidth]{figures/scheduling/schedule_step_01.pdf}
    \caption{Step 1}
  \end{subfigure}\hfill
  \begin{subfigure}[t]{0.24\linewidth}
    \centering
    \includegraphics[width=\linewidth]{figures/scheduling/schedule_step_10.pdf}
    \caption{Step 10}
  \end{subfigure}\hfill
  \begin{subfigure}[t]{0.24\linewidth}
    \centering
    \includegraphics[width=\linewidth]{figures/scheduling/schedule_step_25.pdf}
    \caption{Step 25}
  \end{subfigure}\hfill
  \begin{subfigure}[t]{0.24\linewidth}
    \centering
    \includegraphics[width=\linewidth]{figures/scheduling/schedule_step_50.pdf}
    \caption{Step 50}
  \end{subfigure}
  \caption{Example of a feasible scheduling sequence. The red border marks the
  group being placed at the selected step, while the blue arrow indicates the
  opening side of the plate.}
  \label{fig:sample-schedule}
\end{figure}